\documentclass[12pt,a4paper,oneside]{article}
\usepackage[brazil]{babel}
\usepackage[utf8]{inputenc}
\usepackage[dvipsnames]{xcolor}
\usepackage{listings}
\usepackage{wrapfig}
\usepackage{olymp}

\setcounter{problem}{1}

\begin{document}

\begin{problem}{Quadro na parede}{quadro}{2}

\begingroup
\setlength{\intextsep}{0pt}%
%\setlength{\columnsep}{0pt}

A mãe da Lisa gosta muito de arte e comprou um quadro para colocar na parede da sala. No entanto, ela confundiu as medidas e o quadro não coube na parede.
Felizmente era um quadro de pintura abstrata, então ela pôde simplesmente girar o quadro 90º e ninguém notou.

Como exemplo, suponha que ela tenha comprado um quadro $5 \times 8$ ($5$ metros de altura e $8$ de largura), mas sua parede seja $10 \times 6$. O quadro não cabe na parede porque ele tem $8$ metros de largura e a parede apenas $6$.
Mas, girando o quadro, ele se torna  $8 \times 5$, cabendo tranquilamente na parede $10 \times 6$.

Com receio de errar em novas compras, a mãe da Lisa te pediu um programa para informar as medidas de um quadro e da parede e saber se ele cabe, girado 90º ou não.

% \Notes

\InputFile

A entrada contém duas linhas. A primeira com dois valores inteiros, $A$ e $L$, que são a altura e a largura do quadro. A segunda também tem dois valores inteiros, $X$ e $Y$, que são a altura e a largura da parede.

Restrições: $1 \leq A, L, X, Y \leq 10$.

\endgroup

\OutputFile

Escreva uma linha na saída, com a palavra \texttt{SIM}, se o quadro couber na parede, ou a palavra \texttt{NAO}, caso não caiba, mesmo girado 90º.

% \Notes

\Example

\begin{example}
\exmp{
5 8
10 6
}{
SIM
}%
\end{example}

\begin{example}
\exmp{
8 5
10 6
}{
SIM
}%
\end{example}

\begin{example}
\exmp{
3 5
5 5
}{
SIM
}%
\end{example}

\begin{example}
\exmp{
5 5
4 10
}{
NAO
}%
\end{example}

%Comentários sobre o exemplo, se necessário.

\end{problem}

\end{document}